\documentclass[12pt,a4paper]{article}

% ====== ЯЗЫК И ШРИФТЫ (XeLaTeX) ======
\usepackage{fontspec}
\usepackage{polyglossia}
\setdefaultlanguage{russian}
\setotherlanguage{english}

\setmainfont{CMU Serif}
\setsansfont{CMU Sans Serif}
\setmonofont{CMU Typewriter Text}
\newfontfamily\cyrillicfont{CMU Serif}
\newfontfamily\cyrillicfontsf{CMU Sans Serif}
\newfontfamily\cyrillicfonttt{CMU Typewriter Text}

\usepackage{microtype}


% ====== МАТЕМАТИКА ======
\usepackage{amsmath, amssymb, amsthm, mathtools}

% Русские стандарты математической типографики
\renewcommand{\ge}{\geqslant}
\renewcommand{\le}{\leqslant}

\newcommand{\HH}{\mathrm{H}}
\newcommand{\II}{\mathrm{I}}
\newcommand{\ChiH}{\chi}
\newcommand{\PP}{\mathbb{P}}
\newcommand{\EE}{\mathbb{E}}

% ====== ОФОРМЛЕНИЕ СТРАНИЦЫ ======
\usepackage[left=2cm, right=2cm, top=2.5cm, bottom=2.5cm]{geometry}
\usepackage{parskip} % Красивые отступы между абзацами
\usepackage{xcolor}
\usepackage{enumitem}
\usepackage{array}


% ====== ГИПЕРССЫЛКИ ======
\usepackage[colorlinks=true, linkcolor=defblue!90!black, urlcolor=defblue!90!black, citecolor=probgreen]{hyperref}

% ====== TCOLORBOX (ОКРУЖЕНИЯ ДЛЯ ТЕОРЕМ И ЗАДАЧ) ======
\usepackage[most]{tcolorbox}
\tcbuselibrary{skins, breakable, theorems}

% Цветовая палитра "МЦНМО/ВШЭ"
\definecolor{defblue}{RGB}{20, 75, 140}
\definecolor{thmred}{RGB}{160, 30, 45}
\definecolor{probgreen}{RGB}{25, 110, 55}
\definecolor{warnpurple}{RGB}{130, 30, 130}
\definecolor{insightteal}{RGB}{0, 100, 100}
\definecolor{solgray}{RGB}{90, 90, 90}
\definecolor{codebg}{RGB}{245, 245, 248}

% Определение
\newtcbtheorem[auto counter, number within=section]{definition}{Определение}{
    enhanced, breakable,
    colback=defblue!3, colframe=defblue,
    fonttitle=\bfseries\sffamily, title style={color=defblue},
    attach boxed title to top left={yshift=-2.5mm, xshift=5mm},
    boxed title style={sharp corners},
    boxrule=1pt, drop fuzzy shadow
}{def}

% Теорема
\newtcbtheorem[auto counter, number within=section]{theorem}{Теорема}{
    enhanced, breakable,
    colback=thmred!3, colframe=thmred,
    fonttitle=\bfseries\sffamily, title style={color=thmred},
    attach boxed title to top left={yshift=-2.5mm, xshift=5mm},
    boxed title style={sharp corners},
    boxrule=1pt, drop fuzzy shadow
}{thm}

% Задача
\newtcbtheorem[auto counter, number within=section]{problem}{Задача}{
    enhanced, breakable,
    colback=probgreen!3, colframe=probgreen,
    fonttitle=\bfseries\sffamily, title style={color=probgreen},
    attach boxed title to top left={yshift=-2.5mm, xshift=5mm},
    boxed title style={sharp corners},
    boxrule=1pt, drop fuzzy shadow
}{prob}

% Важное замечание / Парадокс
\newtcolorbox{paradox}[1][]{
    enhanced, breakable,
    colback=warnpurple!4, colframe=warnpurple!90!black,
    title={\textbf{Внимание: #1}},
    fonttitle=\bfseries\sffamily,
    coltitle=black,
    attach title to upper,
    after title={\par\smallskip},
    boxrule=1.5pt, arc=4pt
}

% Олимпиадный инсайт
\newtcolorbox{insight}[1][]{
    enhanced, breakable,
    colback=insightteal!3, colframe=insightteal, leftrule=4pt,
    boxrule=0pt, arc=0pt,
    title={\textbf{Олимпиадный инсайт: #1}},
    coltitle=insightteal, fonttitle=\bfseries\sffamily,
    attach title to upper,
    after title={\par\smallskip}
}

% Решение
\newenvironment{solution}{
    \begin{tcolorbox}[
        enhanced, breakable, blanker,
        borderline west={2.5pt}{0pt}{solgray},
        left=12pt, top=5pt, bottom=5pt,
        before skip=10pt, after skip=15pt
    ]
    \textbf{\textsf{РЕШЕНИЕ.}}\;
}{
    \hfill $\blacksquare$
    \end{tcolorbox}
}



\begin{document}

\begin{center}
    \vspace*{1.0cm}
    {\Huge \bfseries Битовая Вселенная\par}
    \vspace{0.4cm}
    {\Large \itshape Теория информации\par}
    \vspace{0.6cm}
    \rule{0.6\textwidth}{1pt}\par
    \vspace{0.6cm}
    {\large \textbf{Калашников Александр}}\par
    \vspace{0.2cm}
    {\small Wildberries \& Russ, ФКН ВШЭ}
    \vspace{1.0cm}
\end{center}

\tableofcontents
\newpage

\section{Минимум теории, который нужно знать}

\begin{tcolorbox}[colback=codebg, colframe=defblue!80!black, title=\textbf{Договорённости}, fonttitle=\bfseries\sffamily, boxrule=1pt]
\begin{itemize}[leftmargin=1.5em]
    \item Если не сказано иное, все логарифмы --- по основанию $2$, а ответы измеряются в \textbf{битах}.
    \item Если в задаче явно стоит $\ln$, то ответ измеряется в \textbf{натах}.
    \item Главная идея темы: \textbf{Хартли} отвечает за количество информации в худшем случае, \textbf{Шеннон} --- за среднее количество информации при заданных вероятностях.
\end{itemize}
\end{tcolorbox}

\subsection{Информация по Хартли}

\begin{definition}{Информация по Хартли}{hartley}
Если $A$ --- конечное множество равновозможных вариантов, то количество информации в элементе множества равно
$$
\ChiH(A)=\log_2 |A|.
$$
\end{definition}

Если $A \subseteq X \times Y$, то
$$
\ChiH_X(A)=\ChiH(\pi_X(A))=\log_2 |\pi_X(A)|,
\qquad
\ChiH_Y(A)=\ChiH(\pi_Y(A))=\log_2 |\pi_Y(A)|.
$$

\begin{definition}{Условная информация по Хартли}{hartley_cond}
Условная информация в $Y$-компоненте при известном $X$:
$$
\ChiH_{Y|X}(A)=\log_2 \max_{x\in X} |\{y\in Y:(x,y)\in A\}|.
$$
\end{definition}

\begin{theorem}{Неравенства Хартли}{hartley_ineq}
Полезные неравенства:
$$
\ChiH_X(A),\ChiH_Y(A)\le \ChiH(A)\le \ChiH_X(A)+\ChiH_Y(A),
$$
$$
\ChiH(A)\le \ChiH_X(A)+\ChiH_{Y|X}(A).
$$
\end{theorem}

\begin{insight}[Интерпретация Хартли]
\begin{itemize}[leftmargin=1.5em]
    \item Один вопрос с не более чем $m$ исходами несет не более $\log_2 m$ бит.
    \item $k$ таких вопросов различают не более $m^k$ случаев.
    \item Да/нет-вопросы: не более $1$ бита за вопрос.
    \item Взвешивание на чашечных весах: не более $\log_2 3$ бит за взвешивание.
\end{itemize}
\end{insight}

\subsection{Энтропия Шеннона}

\begin{definition}{Энтропия Шеннона}{shannon}
Пусть случайная величина $X$ принимает значения из конечного множества $A$ с вероятностями $p(x)$. Тогда
$$
\HH(X)=-\sum_{x\in A} p(x)\log_2 p(x),
\qquad
0\log_2 0:=0.
$$
\end{definition}

Свойства:
\begin{itemize}[leftmargin=1.5em]
    \item $0\le \HH(X)\le \log_2 |A|$.
    \item $\HH(X)=0$ тогда и только тогда, когда $X$ константа.
    \item $\HH(X)=\log_2 |A|$ тогда и только тогда, когда $X$ равномерна на $A$.
\end{itemize}

Для бинарной случайной величины удобно вводить бинарную энтропию:
$$
h_2(p)=-p\log_2 p-(1-p)\log_2(1-p).
$$

\subsection{Совместная, условная энтропия и взаимная информация}

Для пары случайных величин:
$$
\HH(X,Y)=-\sum_{x,y} p(x,y)\log_2 p(x,y).
$$

\begin{definition}{Условная энтропия}{cond_entropy}
$$
\HH(Y|X)=\sum_x p(x)\HH(Y|X=x)=\HH(X,Y)-\HH(X).
$$
\end{definition}

\begin{theorem}{Цепное правило}{chain}
$$
\HH(X,Y)=\HH(X)+\HH(Y|X)=\HH(Y)+\HH(X|Y).
$$
\end{theorem}

\begin{definition}{Взаимная информация}{mutual}
$$
\II(X:Y)=\HH(X)+\HH(Y)-\HH(X,Y)=\HH(Y)-\HH(Y|X)=\HH(X)-\HH(X|Y).
$$
\end{definition}

\begin{insight}[Ключевые свойства]
\begin{itemize}[leftmargin=1.5em]
    \item $0\le \HH(Y|X)\le \HH(Y)$.
    \item $\II(X:Y)\ge 0$.
    \item $\II(X:Y)=0$ тогда и только тогда, когда $X$ и $Y$ независимы.
    \item Если $Y=f(X)$, то $\HH(Y|X)=0$ и $\HH(X,Y)=\HH(X)$.
    \item Если $X$ и $Y$ независимы, то $\HH(X,Y)=\HH(X)+\HH(Y)$.
    \item Всегда
    $$
    \max\bigl(\HH(X),\HH(Y)\bigr)\le \HH(X,Y)\le \HH(X)+\HH(Y).
    $$
    Верхняя граница достигается при независимости. Нижняя достигается, если одну величину можно сделать функцией другой при сохранении нужных маргиналов.
\end{itemize}
\end{insight}

\subsection{Что особенно важно для олимпиадных задач}

\begin{paradox}[Функциональная зависимость]
\textbf{Если одна величина является функцией другой, то таблицу сопряжения часто можно не строить.}

Например, если $\beta=\lfloor \alpha/2\rfloor$, то
$$
\HH(\alpha,\beta)=\HH(\alpha).
$$
\end{paradox}

\textbf{Если в задаче даны только маргинальные распределения, то:}
\begin{itemize}[leftmargin=1.5em]
    \item максимум совместной энтропии ищется через независимость;
    \item минимум --- через максимально сильную зависимость;
    \item но минимум \textbf{не всегда} равен $\max(\HH(X),\HH(Y))$: надо проверить, можно ли одну величину сделать функцией другой.
\end{itemize}

\textbf{Если задача про классификатор, то информативность выхода о классе --- это взаимная информация.}

Если $Y$ --- истинный класс, а $Z$ --- ответ классификатора, то
$$
\text{информативность}=\II(Y:Z)=\HH(Y)-\HH(Y|Z).
$$

\subsection{Неравенство Фано и классификаторы}

\begin{theorem}{Неравенство Фано}{fano}
Пусть $Y$ принимает $k$ значений, а $\widehat{Y}$ --- предсказание, и
$$
\varepsilon=\PP(\widehat{Y}\ne Y).
$$
Тогда
$$
\HH(Y|\widehat{Y})\le h_2(\varepsilon)+\varepsilon\log_2 (k-1).
$$

Следовательно,
$$
\II(Y:\widehat{Y})\ge \HH(Y)-h_2(\varepsilon)-\varepsilon\log_2 (k-1).
$$
\end{theorem}

Если есть \textbf{отказ} от ответа, то удобно обозначить $Z\in \{1,\dots,k,?\}$, покрытие
$$
q=\PP(Z\ne ?),
$$
а ошибка на покрытых объектах
$$
\varepsilon=\PP(Z\ne Y\mid Z\ne ?).
$$

Тогда
$$
\HH(Y|Z)\le \PP(Z=? )\HH(Y|Z=? )+q\bigl(h_2(\varepsilon)+\varepsilon\log_2 (k-1)\bigr).
$$

Грубая оценка, которая всегда верна:
$$
\HH(Y|Z)\le (1-q)\log_2 k+q\bigl(h_2(\varepsilon)+\varepsilon\log_2 (k-1)\bigr).
$$

Если отказ сам по себе не несет информации о классе, то можно считать
$$
\HH(Y|Z=? )=\HH(Y),
$$
и тогда
$$
\II(Y:Z)\ge q\Bigl(\HH(Y)-h_2(\varepsilon)-\varepsilon\log_2 (k-1)\Bigr).
$$

\newpage
\section{Типовые приёмы решения}

\begin{tcolorbox}[colback=yellow!5, colframe=orange!80!black, title=\textbf{7 золотых правил}, fonttitle=\bfseries]
\begin{enumerate}[leftmargin=1.7em,label=\textbf{\arabic*.}]
    \item \textbf{Сначала поймите, Хартли это или Шеннон.}\\
    Если даны просто варианты без вероятностей --- это Хартли. Если даны вероятности --- это Шеннон.

    \item \textbf{Сначала найдите распределение, потом считайте энтропию.}\\
    В половине задач трудность не в формуле, а в подсчете вероятностей.

    \item \textbf{Проверьте, не является ли одна величина функцией другой.}\\
    Если $Y=f(X)$, то
    $$
    \HH(Y|X)=0,\qquad \HH(X,Y)=\HH(X).
    $$

    \item \textbf{При фиксированных маргиналах сначала подумайте о крайних случаях.}\\
    Максимум --- независимость. Минимум --- максимально сильная зависимость.

    \item \textbf{В задачах про классификаторы переводите все на язык случайных величин.}\\
    Истинный класс --- $Y$, ответ модели --- $Z$, информативность --- $\II(Y:Z)$.

    \item \textbf{Следите за основанием логарифма.}\\
    Если в задаче стоит $\ln$, то ответ надо оставлять через $\ln$, а не переводить в биты.

    \item \textbf{Для пределов полезны суммы рядов:}
    $$
    \sum_{k\ge 1} x^k=\frac{x}{1-x},
    \qquad
    \sum_{k\ge 1} kx^k=\frac{x}{(1-x)^2},
    \qquad |x|<1.
    $$
\end{enumerate}
\end{tcolorbox}

\section{Короткая памятка формул}

\begin{center}
\begin{tabular}{|>{\raggedright\arraybackslash}p{0.35\linewidth}|>{\raggedright\arraybackslash}p{0.57\linewidth}|}
\hline
Хартли & $\ChiH(A)=\log_2|A|$ \\ \hline
Условная информация по Хартли & $\ChiH_{Y|X}(A)=\log_2 \max_{x} |\{y:(x,y)\in A\}|$ \\ \hline
Энтропия Шеннона & $\HH(X)=-\sum p(x)\log_2 p(x)$ \\ \hline
Совместная энтропия & $\HH(X,Y)=\HH(X)+\HH(Y|X)$ \\ \hline
Условная энтропия & $\HH(Y|X)=\HH(X,Y)-\HH(X)$ \\ \hline
Взаимная информация & $\II(X:Y)=\HH(Y)-\HH(Y|X)$ \\ \hline
Независимость & $\HH(X,Y)=\HH(X)+\HH(Y)$ \\ \hline
Функциональная зависимость $Y=f(X)$ & $\HH(Y|X)=0$, $\HH(X,Y)=\HH(X)$ \\ \hline
Фано & $\HH(Y|\widehat Y)\le h_2(\varepsilon)+\varepsilon\log_2(k-1)$ \\ \hline
\end{tabular}
\end{center}

Полезные точные значения:
$$
\HH\Bigl(\frac12,\frac14,\frac14\Bigr)=\frac32,
\qquad
\HH\Bigl(\frac45,\frac1{10},\frac1{10}\Bigr)=\log_2 5-\frac75,
\qquad
\HH\Bigl(\frac25,\frac25,\frac15\Bigr)=\log_2 5-\frac45.
$$

Также полезно помнить:
$$
\log_2 3\approx 1.584963,
\qquad
\log_2 5\approx 2.321928,
\qquad
h_2(0.2)\approx 0.721928.
$$

\newpage
\section{Задачи}

\textbf{Во всех задачах логарифмы по основанию $2$, кроме задачи 9.}

\begin{problem}{Условная информация по Хартли}{hartley_task}
Пусть
$$
X=\{1,2,3\},\qquad Y=\{1,2,3,4,5\},
$$
а
$$
A=\{(1,1),(1,2),(3,3),(1,4),(2,5),(2,2)\}\subseteq X\times Y.
$$
Найдите
$$
\ChiH_{Y|X}(A)=\log_2 \max_{x\in X} |\{y\in Y:(x,y)\in A\}|.
$$
\end{problem}
\begin{solution}
Для $x=1$ имеется $3$ возможных значения $y$, для $x=2$ --- $2$, для $x=3$ --- $1$. Поэтому
$$
\ChiH_{Y|X}(A)=\log_2 3\approx 1.584963.
$$
\end{solution}

\begin{problem}{Адаптивные и неадаптивные вопросы}{binary_search}
Сколько вопросов <<да/нет>> необходимо и достаточно, чтобы однозначно отгадать число из множества $\{1,2,\dots,1000\}$:
\begin{itemize}[leftmargin=1.5em]
    \item адаптивно;
    \item неадаптивно?
\end{itemize}
\end{problem}
\begin{solution}
$$
2^9<1000\le 2^{10},
$$
значит в обоих случаях ответ равен $10$.\\
Адаптивно --- бинарный поиск. Неадаптивно --- можно спрашивать про биты двоичной записи.
\end{solution}

\begin{problem}{Два кубика}{dice}
Бросают два игральных кубика: один четырехгранный с числами $1,2,3,4$, другой шестигранный с числами $1,2,3,4,5,6$. Все исходы равновероятны. Случайная величина $X$ равна модулю разности выпавших чисел. Найдите $\HH(X)$.
\end{problem}
\begin{solution}
Распределение величины $X=|a-b|$:
$$
\PP(X=0)=\frac4{24},\quad
\PP(X=1)=\frac7{24},\quad
\PP(X=2)=\frac6{24},
$$
$$
\PP(X=3)=\frac4{24},\quad
\PP(X=4)=\frac2{24},\quad
\PP(X=5)=\frac1{24}.
$$
Тогда
$$
\HH(X)= -\sum_{k=0}^5 \PP(X=k)\log_2 \PP(X=k)\approx 2.369910.
$$
\end{solution}

\begin{problem}{Функциональная зависимость}{func_dep}
Случайная величина $\alpha$ принимает значения $0,1,2,3$ и имеет распределение
$$
p(0)=\frac18,\qquad p(1)=\frac38,\qquad p(2)=\frac14,\qquad p(3)=\frac14.
$$
Определим
$$
\beta=\left\lfloor \frac{\alpha}{2}\right\rfloor.
$$
Найдите $\HH(\alpha,\beta)$.
\end{problem}
\begin{solution}
Поскольку $\beta$ является функцией от $\alpha$, имеем
$$
\HH(\alpha,\beta)=\HH(\alpha).
$$
Значит
$$
\HH(\alpha,\beta)
=-\frac18\log_2\frac18-\frac38\log_2\frac38-\frac14\log_2\frac14-\frac14\log_2\frac14
=\frac52-\frac38\log_2 3
\approx 1.905639.
$$
\end{solution}

\begin{problem}{Минимум и максимум совместной энтропии}{joint_entropy}
Пусть случайная величина $\alpha$ имеет распределение $\left\{\frac14,\frac14,\frac12\right\}$, а случайная величина $\beta$ имеет распределение $\left\{\frac12,\frac12\right\}$. Найдите минимально возможное и максимально возможное значение $\HH(\alpha,\beta)$.
\end{problem}
\begin{solution}
$$
\HH(\alpha)= -\frac14\log_2\frac14-\frac14\log_2\frac14-\frac12\log_2\frac12=\frac32,
\qquad
\HH(\beta)=1.
$$
Поэтому
$$
\max\bigl(\HH(\alpha),\HH(\beta)\bigr)\le \HH(\alpha,\beta)\le \HH(\alpha)+\HH(\beta).
$$
Нижняя граница достижима: можно сделать $\beta$ функцией от $\alpha$, например выделив состояние $\alpha$ с вероятностью $\frac12$ в одну группу, а два остальных по $\frac14$ --- во вторую. Верхняя достижима при независимости. Значит
$$
\HH_{\min}(\alpha,\beta)=\frac32,
\qquad
\HH_{\max}(\alpha,\beta)=\frac52.
$$
\end{solution}

\begin{problem}{Четность и взаимная информация}{parity_mi}
Пусть $X$ равномерно распределена на множестве $\{0,1,2,3,4,5,6,7\}$, а
$$
Y=X \bmod 2.
$$
Найдите:
$$
\HH(X),\quad \HH(Y),\quad \HH(X|Y),\quad \II(X:Y).
$$
\end{problem}
\begin{solution}
Так как $X$ равномерна на $8$ значениях,
$$
\HH(X)=\log_2 8=3.
$$
Величина $Y$ --- четность, поэтому она равномерна на двух значениях:
$$
\HH(Y)=1.
$$
При известной четности остается $4$ равновозможных значения $X$, значит
$$
\HH(X|Y)=\log_2 4=2.
$$
Наконец,
$$
\II(X:Y)=\HH(X)-\HH(X|Y)=1.
$$
\end{solution}

\begin{problem}{Двоичный симметричный канал}{bsc}
Пусть $X$ --- честный бит, $N\sim \mathrm{Bern}(p)$ независимо от $X$, и
$$
Y=X\oplus N.
$$
Найдите $\HH(Y)$, $\HH(Y|X)$, $\HH(X|Y)$ и $\II(X:Y)$.
\end{problem}
\begin{solution}
Так как $X$ --- честный бит и канал симметричный,
$$
Y \sim \mathrm{Bern}\left(\frac12\right),
\qquad
\HH(Y)=1.
$$
При фиксированном $X$ величина $Y$ совпадает с $X$ с вероятностью $1-p$ и меняется с вероятностью $p$, поэтому
$$
\HH(Y|X)=h_2(p).
$$
По симметрии также
$$
\HH(X|Y)=h_2(p).
$$
Значит
$$
\II(X:Y)=\HH(Y)-\HH(Y|X)=1-h_2(p).
$$
\end{solution}

\begin{problem}{Две монеты}{two_coins}
Подбрасывают две честные монеты. Пусть $X$ --- число выпавших орлов, а $Y$ --- индикатор события <<на первой монете выпал орел>>. Найдите
$$
\HH(X),\quad \HH(Y),\quad \HH(X|Y),\quad \HH(Y|X),\quad \II(X:Y).
$$
\end{problem}
\begin{solution}
У величины $X$ распределение
$$
\PP(X=0)=\frac14,\qquad \PP(X=1)=\frac12,\qquad \PP(X=2)=\frac14,
$$
поэтому
$$
\HH(X)= -\frac14\log_2\frac14-\frac12\log_2\frac12-\frac14\log_2\frac14=\frac32.
$$
Величина $Y$ --- честный бит, значит
$$
\HH(Y)=1.
$$
Если $Y$ известна, то исход второй монеты остается честным битом, значит
$$
\HH(X|Y)=1.
$$
Тогда
$$
\HH(X,Y)=\HH(Y)+\HH(X|Y)=2.
$$
Отсюда
$$
\HH(Y|X)=\HH(X,Y)-\HH(X)=\frac12,
\qquad
\II(X:Y)=\HH(Y)-\HH(Y|X)=\frac12.
$$
\end{solution}

\begin{problem}{Геометрическое распределение}{geom_entropy}
Для вектора вероятностей
$$
\left(\frac23,\frac29,\dots,\frac{2}{3^n},\frac{1}{3^n}\right)
$$
определим
$$
h_n=\sum_i p_i \ln \frac{1}{p_i}.
$$
Найдите
$$
\lim_{n\to\infty} h_n.
$$
\end{problem}
\begin{solution}
Пусть
$$
p_k=\frac{2}{3^k},\qquad k=1,2,\dots,n,
\qquad
p_{n+1}=\frac{1}{3^n}.
$$
Тогда
$$
\lim_{n\to\infty} h_n
=\sum_{k\ge 1}\frac{2}{3^k}\ln\frac{3^k}{2}.
$$
Раскрываем:
$$
\sum_{k\ge 1}\frac{2}{3^k}\ln\frac{3^k}{2}
=
\ln 3\sum_{k\ge 1}k\frac{2}{3^k}
-\ln 2\sum_{k\ge 1}\frac{2}{3^k}.
$$
Используем суммы рядов:
$$
\sum_{k\ge 1}\frac{2}{3^k}=1,
\qquad
\sum_{k\ge 1}k\frac{2}{3^k}=\frac32.
$$
Значит
$$
\lim_{n\to\infty} h_n=\frac32\ln 3-\ln 2.
$$
\end{solution}

\begin{problem}{Классификатор с отказом $(\star)$}{classifier}
На входе классификатора объекты трех классов, причем доли классов относятся как $2:2:1$. Покрытие классификатора равно $0.7$, а точность на покрытых объектах равна $0.8$. Выход классификатора принимает значения из множества $\{1,2,3,?\}$. Найдите минимально возможную информативность выхода классификатора о настоящем классе.
\end{problem}
\begin{solution}
Обозначим истинный класс через $Y$, а ответ классификатора через $Z\in \{1,2,3,?\}$. Нужно минимизировать
$$
\II(Y:Z)=\HH(Y)-\HH(Y|Z).
$$
При отношении классов $2:2:1$ имеем
$$
\HH(Y)=\HH\Bigl(\frac25,\frac25,\frac15\Bigr)=\log_2 5-\frac45.
$$

Оптимально сделать распределение классов на ответе <<$?$>> равномерным:
$$
\PP(Y=i\mid Z=? )=\frac13,\qquad i=1,2,3.
$$
Тогда на покрытой части распределение классов становится
$$
\PP(Y\mid Z\ne ?)=\left(\frac37,\frac37,\frac17\right).
$$
На покрытой части при точности $0.8$ максимум условной энтропии достигается при апостериорных распределениях вида перестановок $(0.8,0.1,0.1)$, поэтому
$$
\HH(Y|Z,Z\ne ?)=\HH\Bigl(\frac45,\frac1{10},\frac1{10}\Bigr)=\log_2 5-\frac75.
$$
Значит
$$
\HH(Y|Z)=0.3\log_2 3+0.7\Bigl(\log_2 5-\frac75\Bigr).
$$
Следовательно,
$$
\II_{\min}(Y:Z)
=\Bigl(\log_2 5-\frac45\Bigr)-0.3\log_2 3-0.7\Bigl(\log_2 5-\frac75\Bigr).
$$
После упрощения:
$$
\II_{\min}(Y:Z)=\frac{3}{10}\log_2\frac53+\frac9{50}\approx 0.401090.
$$
Равенство достижимо.
\end{solution}

\begin{problem}{Минимум совместной энтропии $(\star)$}{joint_min}
Случайная величина $X$ равномерна на трех значениях, а случайная величина $Y$ равномерна на двух значениях. Найдите минимально возможное и максимально возможное значение $\HH(X,Y)$.
\end{problem}
\begin{solution}
Максимум достигается при независимости:
$$
\HH_{\max}(X,Y)=\HH(X)+\HH(Y)=\log_2 3+1\approx 2.584963.
$$

Для минимума надо выбрать максимально зависимое сопряжение. Одно из оптимальных совместных распределений:
$$
\begin{array}{c|cc}
& Y=0 & Y=1\\ \hline
X=1 & \frac13 & 0\\
X=2 & \frac16 & \frac16\\
X=3 & 0 & \frac13
\end{array}
$$
Тогда
$$
\HH(X,Y)=\HH\Bigl(\frac13,\frac16,\frac16,\frac13\Bigr)
=-\frac13\log_2\frac13-\frac16\log_2\frac16-\frac16\log_2\frac16-\frac13\log_2\frac13
=\log_2 3+\frac13.
$$
Поскольку энтропия вогнута по ячейкам таблицы сопряжения, минимум по множеству всех сопряжений достигается в вершине транспортного многогранника; указанная таблица как раз является вершиной. Значит
$$
\HH_{\min}(X,Y)=\log_2 3+\frac13\approx 1.918296.
$$
\end{solution}

\newpage
\section{На что смотреть в первую очередь перед финалом}

\begin{tcolorbox}[colback=codebg, colframe=defblue!80!black, title=\textbf{Чек-лист подготовки}, fonttitle=\bfseries\sffamily, boxrule=1pt]
\begin{itemize}[leftmargin=1.5em]
    \item Уметь \textbf{быстро} считать $\HH(X)$ по распределению.
    \item Мгновенно замечать случаи вида $Y=f(X)$.
    \item Уметь раскладывать
    $$
    \HH(X,Y)=\HH(X)+\HH(Y|X)
    $$
    и
    $$
    \II(X:Y)=\HH(Y)-\HH(Y|X).
    $$
    \item Не путать \textbf{accuracy} и \textbf{coverage}.
    \item Помнить, что в задачах на классификатор информативность выхода --- это взаимная информация, а не просто точность.
    \item Для задач с рядами держать в голове формулы для $\sum x^k$ и $\sum kx^k$.
\end{itemize}
\end{tcolorbox}

\vspace{0.6cm}
\begin{center}
\textit{Итог.}
Для быстрой олимпийской подготовки по этой теме обычно хватает следующего стека:
$$
\text{Хартли} \;\to\; \text{Шеннон} \;\to\; \text{совместная/условная энтропия}
$$
$$
\to\; \text{взаимная информация} \;\to\; \text{Фано и классификаторы}.
$$
\end{center}

\end{document}